% !TEX root = AImath.v5.tex
\chapter*{自序}

在动笔之前，我常被一个朴素而执拗的问题追着走：我们究竟为什么要谈人工智能？是为了追逐新技术的光亮，还是为了在这片光亮背后照见更久远的智性传统——数学、逻辑与经验之间的来回穿梭。写这本书的动机，便是在繁复的术语与迅疾的迭代之上，给读者留一段可以驻足的路径：从问题的提出、到模型的表达、再到工程的落地，层层看清「智能」如何成为可讨论、可计算、可实现的对象。

我坚持从「问题意识」出发。任何技术看似强大，若没有问题意识，也只会迷失在工具的丰盛里。人类的智能之所以动人，是因为它总能把世界重新描述为一组可以被比较、被检验的命题。数学给出形式化的骨架，算法让推理可走，工程则把抽象按到现实的尺度上。三者并非线性递进，而是彼此往复：一个好问题往往逼迫我们改写模型，一个可靠的模型又会反过来修正问题的边界，工程细节最终决定了这些想法能否在真实世界里站住脚。

本书的结构因此尽量保持朴素。每一章都会交代「为何提出这个主题」「我们用怎样的数学去刻画」「具体的算法与实现有哪些关键环节」。读者不必在意术语的完整性，而应关心一个观点能否自洽、能否经受例子与反例的考验。与其把人工智能全盘描述成一架精密仪器，不如把它视为一种持续改进的思维习惯：敢于提出假设，愿意暴露不确定，允许在数据面前修正自信。

在写作上，我尽量避免炫技式的密度。知识固然需要准确，但准确不必与艰涩同义。许多概念其实可以借助朴素比喻来理解：梯度不过是*「朝着更好的方向迈一步」*的量化；正则化不过是*「别让模型把噪声当规律」*的提醒；泛化能力说到底，是*「在陌生数据上仍不失准度」*的朴素期望。我们在复杂与可读之间保持一条窄路，让读者感到「可跟上」，又不会牺牲严谨。

我也希望把「限制」写在书页上。人工智能并不是一切问题的万灵药。从计算理论到工程实践，处处存在代价：算法需要时间与空间，数据需要清洗与标注，硬件有能耗与延迟，更重要的是，真实世界的目标函数往往多维而矛盾。理解这些限制并不让人气馁，反而会使我们在选择上更克制，在评价上更诚实，在改进上更有方向。

读者可能来自不同背景。有工程经验的读者可以从案例出发，逆推背后的数学与假设；偏理论的读者则可以在公式之间找到「可实现性」的锚点；初学者不妨从术语表与插图开始，把大图景先搭起来，再逐步填充细节。我建议以「问题—模型—实现」的节奏来阅读，每到一个关键术语，问自己三个问题：它解决了什么；它假设了什么；它牺牲了什么。

在写作过程中，我反复提醒自己：不要让结论跑在论证前面。数据的可变、环境的复杂、人的目标的多样性，都意味着我们要在多条路径上并行行走。一次实验的成功可能来自巧合；一次算法的失效也许只是样本不够。**把不确定写出来，正是科学精神的底色。**愿这本书在带来启发的同时，也保留谦卑：承认未知，拥抱修正。

如果要用一句话概括本书希望留下的印象，那便是：智能是一种改写世界描述的能力。我们通过数学给出可检验的形式，通过算法让这种形式可运行，再通过工程把它纳入真实约束之中。每一次从「描述」到「实现」的往返，都是对世界的再认识。无论你身处研究、产业还是教育，我都希望你能在这些往返里，找到属于自己的步伐。

最后，感谢一路同行的人：同事与学生提供的讨论，行业伙伴分享的实践，读者反馈出的困惑。它们共同塑造了本书的节奏与取舍。也许我们无法穷尽「智能」的边界，但我们可以在边界附近不断点亮路标。把复杂讲清楚，把清楚做出来，这便是写作与研究最朴素的初心。

愿你在阅读的过程中，不仅收获概念与方法，更收获提出好问题的勇气、承受不确定的耐心以及在现实约束下推进改进的恒心。我们在下一章继续前行，但并不急于抵达；重要的是在每一步停下时，都能回答：为什么在此处停下，下一步又将通向何方。